\section{从成都的票据到河北的堤防}

993--994年成都平原的王小波、李顺起事，1023年益州设交子务，1047--1048年贝州王则起事与黄河商胡决口，构成几种不能相互替代的地方材料问题。史书能确认起事、制度设置和河决的年序；它们却不能把蜀地口号、交子流通、河北的受灾户数都化为同一种“民变”或“财政危机”。交子务说明官府试图接管区域信用，河工与起事则显示地方社会承受的征发、灾害和控制压力各有路径。

这些事件要求把制度文本同地域证据放在一起。交子发行额不等于四川实际货币供给；黄河决口的会要条不等于河北所有聚落的损失；一场起事的传记也不代表全体农户。地方文书、钱币、河道沉积、城址和区域市场材料若能互证，才可追问谁承担运输、借贷、赈济与复原的成本。北宋的财政国家并不是从开封向外均匀铺开，而是在成都、河北、江淮和东南的不同条件中反复碰到边界。

\subsection{票据与河道的两种可见性}

一〇二三年益州交子务接管民间交子并发行官交子，显示区域信用危机被纳入财政工具箱；发行额、准备金和实际流通范围却必须区分制度规定与运行。交子不是开封可直接复制到所有市场的抽象货币，而是在四川商人、铺户、官府和交通条件间取得信用。纸币、钱币、账簿与《宋会要》各自留下不同尺度，不能将一份官办命令或一批出土货币化作区域经济总量。

一〇四八年商胡决口、黄河改道北流，则让河北聚落、耕地、交通与边防同时承担河工和赈济的成本。河道沉积、堤防遗址与《宋会要》可互校改道和维护的主干，受灾人口、河道持续时间和政策成效仍须按地点检验。票据所面对的是信用与流通，河工所面对的是水势、劳役与土地；二者都说明财政行政只能在具体地区的物质条件中发挥作用。
